% !TeX root = ../main.tex
\section{Вывод}

В ходе выполнения лабораторной работы была исследована математическая и имитационная модель тракта цифровой квадратурной демодуляции. По результатам анализа полученных временных диаграмм и спектрограмм можно сделать следующие выводы:

\begin{itemize}
    \item \textbf{Субдискретизация и перенос спектра:} Подтверждена работоспособность метода оцифровки радиосигнала на промежуточной частоте ($210{,}1$~МГц), превышающей частоту дискретизации ($180$~МГц). За счет эффекта алиасинга спектр полезного сигнала корректно зеркально переносится в первую зону Найквиста на частоту $30{,}1$~МГц без наложения и потери информации.
    \item \textbf{Влияние разрядности цифрового гетеродина:} Недостаточная разрядность формирователя опорных колебаний вызывает сильные нелинейные искажения и появление мощных паразитных гармоник, резко сужающих динамический диапазон приемника. Экспериментально подтверждено правило выбора разрядности $n_g = n_a + 2$: при $n_a = 13$~бит и $n_g = 15$~бит шумы квантования ЦГ опускаются ниже уровня шумов АЦП, делая дальнейшее наращивание битности нецелесообразным.
    \item \textbf{Амплитудно-фазовые искажения:} Любое рассогласование параметров квадратурных каналов нарушает их строгую ортогональность, что приводит к неполному подавлению зеркальной спектральной составляющей. Максимальная мощность мешающей компоненты $P_r$ в рамках тестов была зафиксирована при совместном воздействии амплитудной (1{,}5~дБ) и фазовой ($5^\circ$) ошибок, что полностью согласуется с теоретической моделью демодулятора.
\end{itemize}
